
\section*{Общая характеристика работы}

\newcommand{\actuality}{\underline{\textbf{\actualityTXT}}}
\newcommand{\progress}{\underline{\textbf{\progressTXT}}}
\newcommand{\aim}{\underline{{\textbf\aimTXT}}}
\newcommand{\tasks}{\underline{\textbf{\tasksTXT}}}
\newcommand{\novelty}{\underline{\textbf{\noveltyTXT}}}
\newcommand{\influence}{\underline{\textbf{\influenceTXT}}}
\newcommand{\methods}{\underline{\textbf{\methodsTXT}}}
\newcommand{\defpositions}{\underline{\textbf{\defpositionsTXT}}}
\newcommand{\reliability}{\underline{\textbf{\reliabilityTXT}}}
\newcommand{\probation}{\underline{\textbf{\probationTXT}}}
\newcommand{\contribution}{\underline{\textbf{\contributionTXT}}}
\newcommand{\publications}{\underline{\textbf{\publicationsTXT}}}


{\actuality} Предметно-ориентированный язык (англ. Domain-specific language, DSL) -- язык программирования, специализированный для конкретной области применения. Построение такого языка и/или его структура данных отражают специфику решаемых с его помощью задач.

Использование DSL имеет ряд преимуществ:
\begin{itemize}
	\item на этапе проектирования дает возможность создания решений в терминах предметной области, благодаря чему специалисты в данной предметной области могут создавать и модифицировать DSL программы;
	\item решение проблемы с помощью DSL происходит на соответствующем уровне абстракции. Это позволяет экспертам в предметной области понимать и верифицировать DSL-программы;
	\item программы, написанные с использованием DSL лаконичны. Написание DSL с использованием терминов предметной области дает возможность в дальнейшем читать программу достаточно легко;
	\item происходит повышение надежности, эффективности и качества сопровождения. Поскольку на уровне модели операции осуществлять легче, они более эффективны и подвержены меньшему количеству ошибок, чем те же операции на уровне кода;
	\item DSL позволяет на уровне абстракции, соответствующему домену, проводить оптимизацию и валидацию;
	описание домена на одном уровне абстракции можно потом преобразовать в более низкий уровень с подробной детализацией. Таким образом можно дополнять модель на разных этапах разработки.
\end{itemize}

{\aim} Спроектировать язык для описания динамической вычислительной сети.

Для~достижения поставленной цели необходимо было решить следующие {\tasks}:
\begin{enumerate}
  \item Изучить способы реализации DSL.
  \item Составить концептуальную модель языка
  \item Реализовать язык в системе Meta Programming System
  \item Проверить результат на соответствия ограничениям.
\end{enumerate}
    

 % Характеристика работы по структуре во введении и в автореферате не отличается (ГОСТ Р 7.0.11, пункты 5.3.1 и 9.2.1), потому её загружаем из одного и того же внешнего файла, предварительно задав форму выделения некоторым параметрам

%Диссертационная работа была выполнена при поддержке грантов ...

%\underline{\textbf{Объем и структура работы.}} Диссертация состоит из~введения, четырех глав, заключения и~приложения. Полный объем диссертации \textbf{ХХХ}~страниц текста с~\textbf{ХХ}~рисунками и~5~таблицами. Список литературы содержит \textbf{ХХX}~наименование.

%\newpage
\section*{Содержание работы}
Во \underline{\textbf{введении}} обосновывается актуальность исследований, проводимых в рамках данной диссертационной работы, приводится обзор научной литературы по изучаемой проблеме, формулируется цель, ставятся задачи работы, сформулированы научная новизна и практическая значимость представляемой работы.

\underline{\textbf{Первая глава}} посвящена ...

 картинку можно добавить так:
\begin{figure}[ht] 
  \center
  \includegraphics [scale=0.27] {latex}
  \caption{Подпись к картинке.} 
  \label{img:latex}
\end{figure}

Формулы в строку без номера добавляются так:
\[ 
  \lambda_{T_s} = K_x\frac{d{x}}{d{T_s}}, \qquad
  \lambda_{q_s} = K_x\frac{d{x}}{d{q_s}},
\]

\underline{\textbf{Вторая глава}} посвящена исследованию 

\underline{\textbf{Третья глава}} посвящена исследованию 

В \underline{\textbf{четвертой главе}} приведено описание 

В \underline{\textbf{заключении}} приведены основные результаты работы, которые заключаются в следующем:
%% Согласно ГОСТ Р 7.0.11-2011:
%% 5.3.3 В заключении диссертации излагают итоги выполненного исследования, рекомендации, перспективы дальнейшей разработки темы.
%% 9.2.3 В заключении автореферата диссертации излагают итоги данного исследования, рекомендации и перспективы дальнейшей разработки темы.
\begin{enumerate}
	\item Изучены возможности MPI для реализации многопоточности в рамках одного узла.
	\item Спроектирована общая архитектура вычислительного узла ЯГСПП.
	\item Разработана базовая библиотека для реализации ЯГСПП в различных средах.
	\item На основе базовой библиотеки разработана библиотека для реализации ЯГСПП в среде MPI.
\end{enumerate}



%\newpage
При использовании пакета \verb!biblatex! список публикаций автора по теме
диссертации формируется в разделе <<\publications>>\ файла
\verb!../common/characteristic.tex!  при помощи команды \verb!\nocite! 

\ifthenelse{\equal{\thebibliosel}{0}}{% Встроенная реализация с загрузкой файла через движок bibtex8
  \renewcommand{\refname}{\large \authorbibtitle}
  \nocite{*}
  \insertbiblioauthor                          % Подключаем Bib-базы
  %\insertbiblioother   % !!! bibtex не умеет работать с несколькими библиографиями !!!
}{% Реализация пакетом biblatex через движок biber
  \insertbiblioauthor                          % Подключаем Bib-базы
  \insertbiblioother
}

